% ---------------------------------------------------------------------
% 01-preliminaries.tex
% Definitions, standing conventions, and imported tools.
% Keep this section lean: only what later sections actually use.
% ---------------------------------------------------------------------

\section{Preliminaries}
\label{sec:preliminaries}

\subsection{Standing conventions}
\label{sec:conventions}

Record here every convention a reader could reasonably get wrong:
which direction a map is written in, whether rings are commutative,
whether a "smooth" object is assumed second countable, and so on.
A convention stated once and referenced thereafter is better than a
hypothesis repeated in five theorems.

\begin{convention}[Standing assumptions]
\label{conv:standing}
Throughout, let $X$ be a locally compact Hausdorff space and let
$\mu$ be a Radon measure on $X$. All vector spaces are over $\R$
unless stated otherwise.
\end{convention}

\subsection{Definitions}

\begin{definition}[Sample notion]
\label{def:sample}
Call a function $f \colon X \to \R$ \emphterm{admissible} if it is
continuous, $\mu$-integrable, and satisfies
\begin{equation}
\label{eq:admissibility}
  \int_X \abs{f} \dd\mu < \infty
  \quad\text{and}\quad
  \norm{f}_{L^\infty} \le 1 .
\end{equation}
\end{definition}

The two conditions in \cref{eq:admissibility} are independent; neither
implies the other on a general measure space. The following elementary
fact is used repeatedly, so we isolate it.

\begin{lemma}[Uniform bound]
\label{lem:uniform-bound}
Let $f$ be admissible in the sense of \cref{def:sample}. Then for
every $x \in X$,
\begin{equation}
\label{eq:uniform-bound}
  \abs{f(x)} \le \sup_{y \in X} \abs{f(y)} .
\end{equation}
\end{lemma}

\begin{proof}
Immediate from the definition of the supremum. The point of recording
it is that \cref{eq:uniform-bound} is uniform in $x$, which is what
\cref{thm:main} needs.
\end{proof}

\begin{remark}
\label{rem:prelim-remark}
Remarks are numbered in the shared sequence, so \cref{rem:prelim-remark}
can be referenced like any other result. Use remarks for genuine
caveats, not for restating the preceding proof.
\end{remark}
